%%%%%%%%%%%%%%%%%%%%%%%%%%%%%%%%%%%%%%%%%%%%%%%%%%%%%%%%%%%%%%%%%
%Capitulo 2 del cuerpo
%%%%%%%%%%%%%%%%%%%%%%%%%%%%%%%%%%%%%%%%%%%%%%%%%%%%%%%%%%%%%%%%%

\chapter{Nombre del Capítulo tan \\largo que se va a partir controladamente }

\thispagestyle{plain}
\renewcommand{\tablename}{Tabla}

\section{Introducción}

En este capítulo mostraremos las distintas divisiones que tiene el documento, comenzando con el capítulo. Esta sección es la principal del documento y se suele crear un documento por capítulo.

\section{Sección}

La sección es uno de los puntos dentro del capítulo es el segundo en orden jerárquico, justo después del capítulo.
%% forzamos un salto de pagina
\clearpage

\subsection{Subsección}

La subsección es un punto dentro de la sección y ocupa el tercer puesto en el orden jerárquico.

\subsubsection{Subsubsección}

Por ultimo, la subsubsección es el ultimo punto de división posible dentro del documento, ocupa el cuarto puesto en el orden jerárquico, sin embargo no aparecerá con número, tal como vemos, ni aparecerá en el índice de contenidos.
%% forzamos un salto de pagina
\clearpage

\section{Conclusión}

Con esto queda patente las distintas formas de dividir el documento, comenzando con capítulos, y posteriormente creando secciones dentro del mismo.

Si se desea jerarquizar aun más el documento, los títulos de las nuevas partes se deberán introducir a mano y no aparecerán plasmadas en el índice, tal como ocurre con la subsubsección.

\blankpage